% TOP_PROBLEM: {"id": 1200002, "title": "Some Questions — Question 2", "category_id": 17, "category": {"id": 17, "name": "group_theory", "display_name": "Group Theory", "description": "Problems about groups, group actions, representations, and related algebraic structures.", "slug": "group-theory", "order_index": 17, "created_at": "2026-07-31T15:26:25.671Z"}, "status": "open", "problem_number": "AMR-011-0002", "set_id": 13}
% FIRST_POSTED: 2026-10-01
% ENTRY_KIND: statement_only
% UnsolvedMath ID 1200002; source catalogue code AMR-011-0002
% !TeX program = lualatex
% Definition and sources only; this file is not an LLM solution attempt.
\documentclass[11pt,a4paper]{article}
\usepackage[margin=25mm]{geometry}
\usepackage{fontspec}
\setmainfont{lmroman10-regular.otf}[BoldFont=lmroman10-bold.otf,
 ItalicFont=lmroman10-italic.otf,BoldItalicFont=lmroman10-bolditalic.otf]
\usepackage{amsmath,amssymb,mathtools}
\usepackage{unicode-math}
\setmathfont{latinmodern-math.otf}
\AtBeginDocument{\let\setminus\smallsetminus}
\usepackage{xurl,hyperref}
\hypersetup{colorlinks=true,urlcolor=blue}
\setcounter{secnumdepth}{0}
\setlength{\parindent}{0pt}
\setlength{\parskip}{5pt}
\setlength{\emergencystretch}{3em}
\begin{document}
\section*{Some Questions — Question 2}\label{1200002}
{\small UnsolvedMath ID 1200002 · AMR-011-0002 · Group Theory · AMR Open Problem Lists}
\subsection{Definitions and mathematical statement}
Fix a prime $p$. A pro-$p$ group is a compact topological group expressible as an inverse limit of finite $p$-groups. The free pro-$p$ group $F_d^{(p)}$ on $d$ generators is characterized by the universal property that any assignment of its generators to a pro-$p$ group extends uniquely to a continuous homomorphism.

Suppose $G$ has a finite pro-$p$ presentation with $r+2$ generators and $r$ relators for some integer $r\geq0$:
\[
 G=F_{r+2}^{(p)}/\overline{\langle\!\langle R_1,\ldots,R_r
                            \rangle\!\rangle}.
\]
The bar denotes topological closure of the normal subgroup generated by the relators. The numbers refer to an available presentation and are not assumed minimal.

Must there exist an open subgroup $U\leq G$ and a continuous surjective homomorphism $U\to F_2^{(p)}$? Equivalently, one may ask for a continuous quotient which is any nonabelian finitely generated free pro-$p$ group, since such a group maps onto rank two.

Open subgroups have finite index, and all topological homomorphisms and closures are part of the pro-$p$ formulation. The sought conclusion is a quotient of an open subgroup, not merely a free subgroup contained in $G$. Nor is an abstract surjection onto a discrete free group the required object. The prime and presentation may vary, with the subgroup allowed to depend on both. No uniform bound on $[G:U]$ is requested.
\subsection{Short English statement}
Does every pro-p presentation with two more generators than relators have an open subgroup surjecting onto a nonabelian free pro-p group?
\subsection{Sources}
\begin{itemize}
\item[S1] ULAM.AI and the UnsolvedMath Contributors, supplied problems.json, record 1200002 (AMR-011-0002), AMR Open Problem Lists. Catalogue identity and source transcription; CC BY 4.0.
\url{https://huggingface.co/datasets/ulamai/UnsolvedMath}.
\item[S2] Abert - Some questions (pdf) (2010), Question 2. Original source indexed by AMR-011-0002.
\url{https://www.renyi.hu/~abert/questions.pdf}.
\end{itemize}
\end{document}
